%! Unit 2: Quadratic Functions — Factoring Quiz (blank)
% GENERATED BODY: the blank and key bodies are byte-identical except the header title;
% answers live in \slot / work, which the preamble shows or hides.
% One problem per method: GCF, difference of squares, trinomial a=1, trinomial a>1,
% solving by factoring. a > 1 work in the key is SLIP AND SLIDE (as on the example sheet). Front/back: 1--3 on the front, 4--5 on the back; each \vfill
% sits between the key work and the answer line, so the spare room becomes work space. Numbers are distinct from the practice test, unit test and parent visit.
\documentclass[12pt]{article}
\usepackage{algebra2-article}
\usepackage{algebra2-boxes}
\usepackage{grading}
\gradingassignment{alg2-u02-factoring-quiz}
\gradingcourse{algebra_2}
\usepackage{algebra2-graded}
% Question line: forest method label, then the prompt on the same line, over a thin rule
\newcommand{\qhead}[2]{%
  \boxguard[8]\vspace{10pt}\noindent
  {\color{forest}\bfseries #1}\hfill #2\par
  \vspace{-2pt}\noindent{\color{forest}\rule{\linewidth}{0.6pt}}\par}
% THE WORK RULE for a test: \slot{width}{answer} is the ONLY line that differs between
% blank and key, and it differs in the preamble, not the body.
\newcommand{\slot}[2]{\underline{\makebox[#1]{\vphantom{\textbf{#2}}}}}

\begin{document}

\pageheader{Unit 2: Quadratic Functions}{Factoring Quiz}

\vspace{0.06in}
\noindent\textbf{Instructions:} No calculator. Show your work. Factor \textbf{completely}.
Write each final answer in its box. Only what is inside a box is graded. \hfill\textbf{20 pts each}

% ======================================================================
\qhead{1.\ Greatest Common Factor}{Factor completely: \quad $6x^3-15x^2+12x$}
\begin{work}
  \text{GCF} &= 3x \\
  6x^3-15x^2+12x &= 3x(2x^2)-3x(5x)+3x(4) \\
  &= 3x(2x^2-5x+4)
\end{work}
\vfill
\noindent\hfill Answer: \inblank{\slot{4.2cm}{$3x(2x^2-5x+4)$}}{\gradedblank[id=q1, type=short, points=20, answer={3x(2x^2-5x+4)}, check=expression]{4.2cm-2pt}{18pt}}

% ======================================================================
\qhead{2.\ Difference of Squares}{Factor completely: \quad $9x^2-49$}
\begin{work}
  9x^2-49 &= (3x)^2-(7)^2 \\
  a^2-b^2 &= (a+b)(a-b)
\end{work}
\vfill
\noindent\hfill Answer: \inblank{\slot{4.2cm}{$(3x+7)(3x-7)$}}{\gradedblank[id=q2, type=short, points=20, answer={(3x+7)(3x-7)}, check=expression]{4.2cm-2pt}{18pt}}

% ======================================================================
\qhead{3.\ Trinomial, $a=1$}{Factor completely: \quad $x^2+3x-40$}
\begin{work}
  \text{multiply to } {-40},\ \text{add to } 3 &: \quad 8 \text{ and } {-5} \\
  x^2+3x-40 &= (x+8)(x-5)
\end{work}
\vfill
\noindent\hfill Answer: \inblank{\slot{4.2cm}{$(x+8)(x-5)$}}{\gradedblank[id=q3, type=short, points=20, answer={(x+8)(x-5)}, check=expression]{4.2cm-2pt}{18pt}}

% ======================================================================
\newpage
\qhead{4.\ Trinomial, $a>1$}{Factor completely: \quad $3x^2-10x-8$}
\begin{work}
  \text{Slip: } & x^2-10x-24 \\
  \text{Factor: } & (x-12)(x+2) \\
  \text{Slide: } & \left(x-\tfrac{12}{3}\right)\left(x+\tfrac{2}{3}\right) \\
  \text{Simplify: } & (x-4)\left(x+\tfrac{2}{3}\right) \\
  \text{Slide up: } & (x-4)(3x+2)
\end{work}
\vfill
\noindent\hfill Answer: \inblank{\slot{4.2cm}{$(x-4)(3x+2)$}}{\gradedblank[id=q4, type=short, points=20, answer={(x-4)(3x+2)}, check=expression]{4.2cm-2pt}{18pt}}

% ======================================================================
\qhead{5.\ Solving a Quadratic Equation by Factoring}{Solve by factoring: \quad $2x^2+5x=12$}
\begin{work}
  & 2x^2+5x-12 = 0 \\
  \text{Slip: } & x^2+5x-24 \\
  \text{Factor: } & (x+8)(x-3) \\
  \text{Slide: } & \left(x+\tfrac{8}{2}\right)\left(x-\tfrac{3}{2}\right) \\
  \text{Slide up: } & (x+4)(2x-3) = 0 \\
  & x+4=0 \quad\text{or}\quad 2x-3=0
\end{work}
\vfill
\noindent\hfill Solutions: \inblank{\slot{4.2cm}{$x=-4$ or $x=\tfrac{3}{2}$}}{\gradedblank[id=q5, type=short, points=20, answer={-4, 3/2}, check=set, display={$x=-4$ or $x=3/2$}]{4.2cm-2pt}{18pt}}

\end{document}
